% Introduction (PUB-REF-0001 AN-21 F-BOR).  Sources: ADR-0007 DEC-05 (the
% RTL boundary); docs/pubs/originals/borland-turbo-pascal.md s1, s3, s4.3.
\Gates{test-pubs-build, test-pubs-tps}
\IntroductionChapter

This book describes the standard procedures and functions of Turbo
Initech\reg{} in alphabetical order.  Each entry tells you what the routine
does, gives its declaration, explains its parameters and its behavior in
every case, and shows a complete program that uses it, with the output the
program produces.  Appendix A describes the entry points of the run-time
library for programmers who read the assembly language the compiler
writes.

\section{What is in the library}
\ix{run-time library}
The run-time library of Turbo Initech is deliberately small.  It provides
output to the serial port, conversion between characters, Boolean values
and integers, the length of a string, and five procedures that read and
write files in blocks of bytes through InitechDOS.  Everything else a
program needs, it writes itself in Pascal: the compiler's own conversion of
integers to decimal text, for example, is an ordinary Turbo Initech
function built from \rw{div}, \rw{mod}, \std{chr}{Chr} and string
concatenation.

\section{The Turbo Initech manuals}
\ix{manuals}
\SideHead{User's Guide}The \doctitle{Turbo Initech User's Guide}
(1077-0000-10) shows you how to compile and run programs and teaches the
language by example.

\SideHead{Programmer's Guide}The \doctitle{Turbo Initech Programmer's
Guide} (1077-0000-11) is the definition of the language.  Its Chapter 9
lists the standard routines by purpose.

\SideHead{Library Reference}This book (1077-0000-12) describes each
standard routine.

\SideHead{Self-Compilation Guide}The \doctitle{Turbo Initech
Self-Compilation Guide} (1077-0000-13) describes how Turbo Initech is built
from its own source.

\SideHead{Quick Reference}The \doctitle{Turbo Initech Quick Reference}
(1077-0000-20) summarizes the syntax, the routines and the messages.

\section{Typefaces used in these books}
\ix{typefaces}
\begin{BorTable}{Typefaces}{p{38mm}p{62mm}}
\Hd{Example} & \Hd{Meaning}\EndHd
\texttt{Writeln(Total);} & Program text, declarations, and what programs
  print\\
\texttt{\bfseries begin} & A reserved word inside program text\\
\textbf{begin} & A reserved word in running text\\
\textit{Count} & A parameter, an identifier of a program, or a routine
  named in running text\\
\MarginFace Margin note & A note on the paragraph beside it
\end{BorTable}

The arrow in the margin marks a remark that you should not overlook.

\section{How to contact Initech}
\ix{technical support}
If you have a question that this book does not answer, call the Initech
Help Desk at 555-0125, extension 2504, or write to \InitechAddressLine.
